\section{Model Evaluation}

\subsection{Advantages}

The developed model system is systematic and forward-looking, characterized by the following three strengths:

\begin{enumerate}
    \item \textbf{Economic and Logistics Optimization:} The MY-HCTO-PM model transcends traditional single-objective optimization by integrating a ``two-factor cost structure'' and Wright's Law. This framework captures the physical trajectory of declining space transportation costs, avoiding the distortions of purely mathematical forecasting. Furthermore, Pareto Mapping provides practical decision-making trade-offs between mission cycles and budgetary constraints.
    \item \textbf{High-Precision Resource Management:} The water demand model for a lunar colony of 100,000 inhabitants accounts for complex fluxes—including physiological metabolism, ecological agriculture, industrial production, and technical losses such as airlock leakage. Validated against International Space Station (ISS) and closed-ecosystem experimental data, this multi-dimensional balance ensures a robust physical foundation for Earth-Moon logistics.
    \item \textbf{Innovative Sustainability Framework:} By incorporating Life Cycle Assessment (LCA) and stratospheric aerosol carrying capacity, the study introduces ``environmental cost'' as a third dimension to the time-cost trade-off system. This elevates the framework to a sustainable development model, quantitatively demonstrating the unique ecological advantages of space elevators for ultra-large-scale transportation.
\end{enumerate}

\subsection{Limitations}

Despite its performance, the model possesses certain limitations:

\begin{enumerate}
    \item \textbf{Linearity in Technology Iteration:} The model relies on historical parameters (e.g., 5\% annual decay), which may fail to capture the nonlinear, disruptive breakthroughs characteristic of aerospace technology. Sudden advancements in materials or propulsion (``black swan'' events) are difficult to account for within the current exponential decay framework.
    \item \textbf{Steady-State Assumptions:} The model simplifies the nonlinear coupling between lunar population growth and infrastructure expansion. By assuming a steady state post-construction, it may overlook the dynamic complexities of population migration and industrial ramping, potentially introducing biases during early expansion phases.
\end{enumerate}